\begin{abstract}
本文围绕休闲零食连锁店商品销量预测问题，基于 2020 年 4 月至 2022 年 3 月的销售明细，以及天气、节假日和促销活动等外部信息，构造日期--门店--商品粒度的日销量面板。数据预处理阶段区分正向销售、净销量和负向调整，将正向销量作为主要预测目标，并在门店、商品、类别和门店--商品组合四个层级开展建模分析。

针对题目四问，本文采用由简单模型到综合模型的分层路线。问题一比较移动平均、同星期均值和简单指数平滑，建立门店、商品和门店--商品组合的基础预测模型；问题二以题目给定类别作为同类零食的聚合依据，并用相关性分析刻画商品销量的同步波动；问题三采用合并天气、对数销量和历史控制的固定效应回归，分析外部变量与销量之间的条件统计关联；问题四将历史销量、类别、日历和外部变量纳入 Ridge 综合模型，并采用严格 7 日递推验证评价一次性预测未来 7 天的效果。

结果表明，简单、可解释的基准模型具有很强竞争力。按更贴近运营汇总的 7 日窗口总量验证，门店层和商品层 WAPE 分别约为 22\% 和 21\%；问题二日滚动一步验证下类别聚合后直接预测 WAPE=34.74\%。门店--商品细粒度层面因大量间歇和零销量序列放大 WAPE 分母而误差偏高。问题三主回归样本量为 58517，调整 $R^2=0.370$，活动日与对数销量呈正向统计关联，但外部因素结论只能解释为观察数据下的条件关联。问题四严格 7 日递推验证中，问题一门店--商品简单指数平滑 WAPE=76.35\%，综合 Ridge WAPE=76.68\%；低销量序列采用指数平滑、常规序列采用综合 Ridge 的混合策略 WAPE=75.85\%，相对简单指数平滑降低 0.49 个百分点，但配对检验和 bootstrap 置信区间显示该差异未达统计显著。因此，综合模型的价值主要在于统一融合历史、日历、类别和外部变量，支撑外部情景下的条件预测与因素解释，而非证明细粒度绝对误差显著下降。
\end{abstract}

\noindent\textbf{关键词：}销量预测；指数平滑；类别聚合；固定效应回归；严格递推验证；WAPE
